\section{Periodic tails are fixed points}
\label{sec:periodic-fixed-points}

The matrix language answers a question that the nested notation conceals:
what kind of number can a repeating continued fraction represent?
Throughout this discussion the continued fractions are \emph{simple}:
$a_0\in\Z$ and $a_n\ge1$ for $n\ge1$.

\begin{definition}
An infinite continued fraction is \emph{eventually periodic} if its digits
repeat after a finite prefix. It is \emph{purely periodic} if repetition
starts at the first digit:
\[
 [\overline{b_0,\ldots,b_{\ell-1}}]
 =[b_0;b_1,\ldots,b_{\ell-1},b_0,b_1,\ldots].
\]
In a purely periodic simple continued fraction, every $b_i$ is positive,
including $b_0$. Rational numbers have finite simple continued fractions;
we do not call those eventually periodic here.
\end{definition}

\begin{definition}
A \emph{quadratic irrational} (or quadratic surd) is an irrational real
root of a quadratic polynomial with integer coefficients. Its other root
is its \emph{conjugate}, written $x'$. Equivalently,
\[
 x=u+v\sqrt D,\qquad x'=u-v\sqrt D
 \quad (u,v\in\Q,\ v\ne0),
\]
where $D$ is a positive nonsquare integer. Replacing $\sqrt D$ by
$-\sqrt D$ respects addition, multiplication, and division by nonzero
numbers. In particular, $(T_M(x))'=T_M(x')$ for rational $M$.
\end{definition}

\begin{lemma}[A nontrivial fixed point is quadratic]
\label{lem:mobius-fixed-point}
If an invertible rational matrix
$M=\begin{psmallmatrix}a&b\\c&d\end{psmallmatrix}$ is not a scalar
multiple of the identity and fixes an irrational number $t$ under its
linear fractional action, then $t$ is a quadratic irrational.
\end{lemma}
\begin{proof}
Multiplying $t=(at+b)/(ct+d)$ by its denominator gives
\begin{equation}\label{eq:mobius-fixed-point}
 ct^2+(d-a)t-b=0.
\end{equation}
The coefficients are not all zero: that would say $c=b=0$ and $d=a$,
making $M$ scalar. Moreover, $c$ cannot be zero, since a nonzero linear
polynomial with rational coefficients cannot vanish at an irrational
number. Thus this is a genuine quadratic equation. Clear denominators to
obtain integer coefficients.
\end{proof}

\begin{warningbox}
The qualification \emph{nontrivial} matters. The identity transformation
fixes every real number, so $t=T_I(t)$ by itself says nothing about $t$.
For a period matrix there is no such problem: its lower-left entry is
strictly positive.
\end{warningbox}

\begin{theorem}[Lagrange]
\label{thm:lagrange-periodicity}
A real number has an eventually periodic simple continued fraction if and
only if it is a quadratic irrational.
\end{theorem}

\begin{proof}[Proof that periodicity implies a quadratic equation]
First let $x=[\overline{b_0,\ldots,b_{\ell-1}}]$. Removing one complete
period leaves the same infinite tail, so
\[
 x=T_N(x),\qquad
 N=D(b_0)\cdots D(b_{\ell-1})
   =\begin{pmatrix}A&B\\C&D_0\end{pmatrix}.
\]
Because every $b_i\ge1$, we have $B,C>0$. Also
$\det N=(-1)^\ell$. The fixed-point equation is
\[
 Cx^2+(D_0-A)x-B=0.
\]
The alternating-convergent argument also applies to an infinite
positive digit sequence: it gives a limit and an error at most
$1/(q_kq_{k+1})$. That limit cannot be rational. Indeed, the
convergents are in lowest terms and $q_k\to\infty$, so they are
eventually distinct from any fixed rational $p/q$. Their distance
from $p/q$ is then at least $1/(q q_k)$, contradicting the error
bound once $q_{k+1}>q$. Thus $x$ is quadratic.

Now allow a prefix:
\[
 y=[a_0;\ldots,a_{r-1},\overline{b_0,\ldots,b_{\ell-1}}]
   =T_M(x),\qquad M=D(a_0)\cdots D(a_{r-1}).
\]
If the prefix is empty, take $M=I$. Since $T_M$ is invertible,
substituting $x=T_{M^{-1}}(y)$ into $x=T_N(x)$ gives
\begin{equation}\label{eq:prefix-conjugation}
 \boxed{\ y=T_{MNM^{-1}}(y).\ }
\end{equation}
Both $M$ and $N$ have determinant $\pm1$, so $K=MNM^{-1}$ is an
integer matrix. It is not scalar: if $K=\lambda I$, then
$N=M^{-1}KM=\lambda I$, contrary to $C>0$.

Finally, write $M=\begin{psmallmatrix}a&b\\c&d\end{psmallmatrix}$.
The number $y$ is irrational. Otherwise the inverse formula
$x=T_{M^{-1}}(y)$ would make $x$ rational. Its denominator cannot vanish
at $y=T_M(x)$: direct substitution gives
$a-cy=(ad-bc)/(cx+d)\ne0$.
The fixed-point lemma applied to $K$ therefore gives a genuine quadratic
equation for $y$ as well.
\end{proof}

\begin{bigidea}
A repeating tail is a fixed point. Attaching a finite prefix
\emph{conjugates the transformation that fixes it}. This is the missing
algebraic step between a purely periodic tail and an eventually periodic
number; it requires only matrix multiplication and inversion.
\end{bigidea}

\begin{example}[A prefix changes the equation, not its degree]
Take
\[
 x=[\overline{1,2}],\qquad
 N=D(1)D(2)=\begin{pmatrix}3&1\\2&1\end{pmatrix}.
\]
Then $2x^2-2x-1=0$, so $x=(1+\sqrt3)/2$. Prefix the digit $1$:
\[
 y=[1;\overline{1,2}]=1+\frac1x,\qquad
 M=D(1)=\begin{pmatrix}1&1\\1&0\end{pmatrix}.
\]
The conjugated matrix is
\[
 MNM^{-1}
 =\begin{pmatrix}1&1\\1&0\end{pmatrix}
  \begin{pmatrix}3&1\\2&1\end{pmatrix}
  \begin{pmatrix}0&1\\1&-1\end{pmatrix}
 =\begin{pmatrix}2&3\\1&2\end{pmatrix}.
\]
Thus $y=(2y+3)/(y+2)$, giving $y^2=3$. Since $y>0$, $y=\sqrt3$.
The same matrix will reappear as a generator of Pell solutions.
\end{example}

\section{Why every quadratic irrational repeats}
\label{sec:quadratic-finite-states}

We now prove the other direction of Lagrange's theorem. There are two
tasks: find a bounded set of integer states, and show that the continued
fraction eventually enters it. The conjugate supplies exactly the
inequalities we need.

\begin{definition}
A quadratic irrational $x$ is \emph{reduced} if
\[
 x>1\qquad\text{and}\qquad -1<x'<0.
\]
\end{definition}

\subsection*{Step 1: every quadratic irrational has a reduced tail}
Write $x_0=x$, $a_n=\floor{x_n}$, and
$x_{n+1}=1/(x_n-a_n)$. For $n\ge1$, $x_n>1$ and $a_n\ge1$.
Conjugating the recurrence gives
\[
 x_{n+1}'=\frac1{x_n'-a_n}.
\]
Consequently, once $x_n'<0$ with $n\ge1$, the next conjugate belongs to
$(-1,0)$. Every later complete quotient is reduced too.

Why must a negative conjugate appear? Invert the matrix for the first
$n$ digits and then conjugate. For $n\ge2$ this gives
\[
 x_n'=-\frac{q_{n-2}x'-p_{n-2}}{q_{n-1}x'-p_{n-1}}.
\]
The convergents $p_k/q_k$ tend to $x$, and $x'\ne x$. Hence, for all
sufficiently large $k$, the numbers
$x'-p_k/q_k$ have the same sign as $x'-x$.
The numerator and denominator in the displayed ratio therefore have
the same sign for large $n$, since the $q_k$ are positive.
It follows that $x_n'<0$ eventually, as required.
This argument also applies when the original $x$ is negative or less
than $1$.

\subsection*{Step 2: integer states evolve by an exact recurrence}
Choose an equation $Ax^2+Bx+C=0$ with integer coefficients and $A\ne0$.
After changing all three signs if necessary, the quadratic formula reads
\[
 x=\frac{-B+\sqrt D}{2A},
 \qquad D=B^2-4AC>0.
\]
Here $D$ is a nonsquare integer; it need not be squarefree. Put
$P_0=-B$ and $Q_0=2A$. Then $Q_0\ne0$ and
$Q_0\divides D-P_0^2$, because $D-P_0^2=-4AC$.
The sign change above is important: $Q_0$ may initially be negative.

Suppose $x_n=(P_n+\sqrt D)/Q_n$ with integers $P_n,Q_n$ satisfying
$Q_n\divides D-P_n^2$. Rationalizing $1/(x_n-a_n)$ yields
\begin{equation}\label{eq:quadratic-state-recurrence}
 \boxed{\quad
 P_{n+1}=a_nQ_n-P_n,\qquad
 Q_{n+1}=\frac{D-P_{n+1}^2}{Q_n}.
 \quad}
\end{equation}
Indeed,
\[
 \frac1{x_n-a_n}
 =\frac{Q_n}{\sqrt D-P_{n+1}}
 =\frac{P_{n+1}+\sqrt D}{(D-P_{n+1}^2)/Q_n}.
\]
Since $P_{n+1}\congr-P_n\pmod{Q_n}$, the new $Q_{n+1}$ is an integer.
It is nonzero because $D$ is not a square. The equality
$D-P_{n+1}^2=Q_nQ_{n+1}$ supplies the divisibility condition for the
next step. Thus the representation persists forever with the same $D$.

\subsection*{Step 3: reduced states lie in a finite box}
For a reduced complete quotient,
\[
 x_n=\frac{P_n+\sqrt D}{Q_n}>1,\qquad
 -1<\frac{P_n-\sqrt D}{Q_n}<0.
\]
Their difference is $2\sqrt D/Q_n>1$, so
$0<Q_n<2\sqrt D$. The conjugate inequalities then give
$\sqrt D-Q_n<P_n<\sqrt D$. In particular,
\begin{equation}\label{eq:reduced-state-bounds}
 \boxed{\quad |P_n|<\sqrt D,\qquad 0<Q_n<2\sqrt D.\quad}
\end{equation}
Only finitely many integer pairs satisfy these bounds.
By Step 1 all sufficiently late states lie in this box, and by Step 2
each state determines its next state, including the digit
$a_n=\floor{(P_n+\sqrt D)/Q_n}$.
Two states must repeat. Every later state, and hence every later digit,
then repeats with the same period. This completes the proof of
Theorem~\ref{thm:lagrange-periodicity}.

\begin{example}[The states for $\sqrt3$]
Starting with $(P_0,Q_0)=(0,1)$ and $D=3$ gives
\begin{center}
\begin{tabular}{ccccc}
\toprule
$n$ & $P_n$ & $Q_n$ & $x_n$ & $a_n$\\
\midrule
$0$ & $0$ & $1$ & $\sqrt3$ & $1$\\
$1$ & $1$ & $2$ & $(1+\sqrt3)/2$ & $1$\\
$2$ & $1$ & $1$ & $1+\sqrt3$ & $2$\\
$3$ & $1$ & $2$ & $(1+\sqrt3)/2$ & $1$\\
\bottomrule
\end{tabular}
\end{center}
The initial state is not reduced. The next two are reduced and form a
cycle, giving $\sqrt3=[1;\overline{1,2}]$.
\end{example}

\section{Galois: when repetition starts immediately}
\label{sec:galois}

The fixed-point equation tells us that periodic continued fractions are
quadratic. It does \emph{not} tell us which quadratic irrationals have
no initial nonrepeating segment. That distinction is Galois's theorem.

\begin{theorem}[Galois's pure-periodicity criterion]
\label{thm:galois}
The simple continued fraction of a quadratic irrational $x$ is purely
periodic if and only if $x$ is reduced:
\[
 x>1\qquad\text{and}\qquad -1<x'<0.
\]
\end{theorem}
\begin{proof}
Suppose first that $x$ is purely periodic with period length $\ell$.
The fixed-point argument shows that it is quadratic, and Step 1 above
shows that all sufficiently late complete quotients are reduced.
But $x_{k\ell}=x_0=x$ for every $k$. Taking $k$ large proves that
$x$ itself is reduced.

Conversely, suppose $x$ is reduced. All its complete quotients remain
reduced, and the finite-state proof gives $x_i=x_j$ for some $i<j$.
To conclude that the period starts at $0$, we need one more fact:
\emph{a reduced complete quotient has at most one reduced predecessor}.

Indeed, suppose the continued-fraction step sends a reduced $u$ to $z$,
and let $a=\floor u$. Then $z=1/(u-a)$, so
\[
 -\frac1{z'}=a-u'.
\]
Since $-1<u'<0$, we have $a<-1/z'<a+1$. Therefore
\begin{equation}\label{eq:reduced-predecessor}
 \boxed{\quad
 a=\floor{-1/z'},\qquad u=a+\frac1z.
 \quad}
\end{equation}
Both the previous digit and the previous reduced number are determined
by $z$. If $i>0$, equality $x_i=x_j$ therefore forces
$x_{i-1}=x_{j-1}$. Repeating the argument gives $x_0=x_{j-i}$.
The deterministic continued-fraction algorithm now repeats from its
first step, proving pure periodicity.
\end{proof}

\begin{bigidea}
Finitely many states give \emph{eventual} periodicity.
Recovering the previous reduced state removes the possible initial
segment and gives \emph{pure} periodicity. These are different steps.
For instance, the finite map $0\mapsto1$, $1\mapsto1$ has an initial
segment when started at $0$; finiteness alone does not exclude one.
\end{bigidea}

\begin{example}[Two inequalities, both necessary]
The number $\sqrt2=[1;\overline2]$ is not reduced, because
$-\sqrt2<-1$. It is eventually but not purely periodic.
In contrast,
\[
 1+\sqrt2=[\overline2],
 \qquad \frac{1+\sqrt5}{2}=[\overline1]
\]
are reduced and purely periodic.
The number $\sqrt2/2$ has conjugate in $(-1,0)$ but is less than $1$,
so it is not purely periodic either.

A displayed prefix does not necessarily represent a genuine preperiod:
\[
 [2;\overline{1,2}]=1+\sqrt3=[\overline{2,1}].
\]
Here the conjugate $1-\sqrt3$ lies in $(-1,0)$, so Galois correctly
detects pure periodicity despite the first way of writing the digits.
\end{example}

\begin{corollary}[The conjugate reverses the period]
If $x=[\overline{b_0,\ldots,b_{\ell-1}}]$, then
\[
 -\frac1{x'}=[\overline{b_{\ell-1},\ldots,b_0}].
\]
\end{corollary}
\begin{proof}
Each $D(b_i)$ is symmetric, so reversing the digits transposes their
product $N=\begin{psmallmatrix}A&B\\C&D_0\end{psmallmatrix}$.
Its fixed-point equation is $Cx^2+(D_0-A)x-B=0$. Substituting
$x'=-1/z$ gives
\[
 Bz^2+(D_0-A)z-C=0,
\]
which is the fixed-point equation for $N^T$. Since $B,C>0$, this
quadratic has exactly one positive root. Both $-1/x'$ and the
reversed periodic continued fraction are that root.
\end{proof}

\begin{checkpoint}[title={Three names, three different conclusions}]
\textbf{Legendre}: an exceptionally close rational approximation must
be a convergent.
\quad\textbf{Lagrange}: the numbers with eventually periodic simple
continued fractions are exactly the quadratic irrationals.
\quad\textbf{Galois}: among them, pure periodicity means exactly
$x>1$ and $-1<x'<0$.
The first direction of Lagrange is the fixed-point calculation;
Galois's additional content is the precise conjugate criterion.
\end{checkpoint}
